
\sgaStableTarget{sga6.e10.eq.1-1-3}@@M0@@
\uline{where} @@M1@@.

By immediate induction on @@M2@@, one obtains:

\sgaStableTarget{sga6.e10.cor.1-1-4}\textbf{Corollary 1.1.4.} \uline{For every noetherian scheme} @@M3@@ \uline{and every integer} @@M4@@, @@M5@@ \uline{is the subgroup of} @@M6@@ \uline{generated by the classes} @@M7@@\uline{, where} @@M8@@ \uline{runs through the set of integral closed subschemes of dimension} @@M9@@ \uline{in} @@M10@@.

\emph{Proof of} 1.1.2: Let @@M11@@. Since @@M12@@ is a noetherian space of dimension @@M13@@, the elements of @@M14@@ are maximal in @@M15@@, hence finite in number, and @@M16@@ is finite for every @@M17@@. This proves the first assertion.

Let @@M18@@ be the category of coherent @@M19@@-modules, and let @@M20@@ be the subset of @@M21@@ consisting of the coherent @@M22@@-modules concentrated on @@M23@@ that satisfy condition (1.1.3). Clearly @@M24@@ contains the sheaves @@M25@@ with @@M26@@. Thus it is enough to show that @@M27@@ is an ``exact'' subset of @@M28@@, since the dévissage lemma (EGA III 3.1.2) then gives @@M29@@.

Thus let
\sgaStableTarget{sga6.e10.eq.1-1-5}@@M30@@
be an exact sequence of @@M31@@-modules with two of its terms in @@M32@@. Then the sheaves @@M33@@, @@M34@@ and @@M35@@ are all concentrated on @@M36@@, so the lengths @@M37@@, @@M38@@ and @@M39@@ of the @@M40@@-modules @@M41@@, @@M42@@ and @@M43@@ are finite

